\section{A prioritized research program}
The immediate objective is a structural account of integral Springer--Borel images. The Cayley results form a substantial related project, but reopening a long exceptional census would not advance the current ordinary-image question. The following priorities separate an active proof problem, consolidation needed for publication, and longer-term extensions.

\subsection{Priority one: a short presentation of classical saturations}
\textbf{Target.} Prove or refute $J_e=K_e$ for $\ell\ge3$ and $1\le e<\ell$ in the square-zero defect model. Together with the proved denominator formula, this would turn the complete B/D square Smith theory into a uniform small-generator description of the saturated ring.

\textbf{Route.} Work with the coefficient images of $\bigwedge^\ell(1+\sum_{j\ge c}a_jT^j)$ over $\Z$. Seek a straightening order with monic reductions and a basis whose degree distribution matches the defect variety. The two-variable proof supplies the correct standard: integral freeness plus the known rank, not rational rank alone. An alternative route is an integral incidence filtration that identifies these operator images with primitive geometric subquotients.

\textbf{Decisive test.} The argument must retain all coefficient operators, divided-power multiplicities, and boundary terms. A proof after inverting factorials is partial. A counterexample must identify a specific integral class in $K_e/J_e$ and explain whether the failure is rank or torsion. Either outcome would be progress: the former proves the proposed presentation; the latter identifies the missing generator mechanism.

\textbf{Literature comparison.} Compare the exact Grassmannian quotient and integral coefficient ring with truncated symmetric ideals, Tanisaki-type presentations, and integral current-algebra modules. Similar Hilbert series are useful clues but not identifications. In particular stable-plane conditions and annihilated-flag conditions differ, and characteristic-zero module bases cannot be imported as integral primitivity theorems without a compatible lattice map.

\subsection{Priority two: consolidate the ordinary theorem package}
\textbf{Target.} Prepare the source-lattice theorem, B/D Smith formulas, adapted bases, and geometric denominator theorem for expert review as a coherent package. A review should separate published integral Schubert inputs from the additional primitive-prefix and gluing arguments.

\textbf{Route.} Supply a minimal dependency graph, fix a single notation dictionary, and expose the key integral assertions as standalone lemmas. The full D gluing step and the use of Chow envelopes deserve particular attention. Replace any computationally observed statement with a proof or restrict its claim to the certified cases. The new statement is a guide to this work, not a substitute for that review.

\textbf{Success.} A reader can verify the theorem without reconstructing a conversation or traversing historical checkpoints. The statement should explain exactly how the new finite integral results refine Parts II and III, while identifying stable injectivity and the underlying Chern expressions as published inputs. A separate novelty comparison should identify the exact integral-image formulation, rather than claim novelty for the Gaussian counting polynomial or standard straightening machinery.

\subsection{Priority three: complete a coherent specialization range}
\textbf{Target.} Determine the complete type A projective K-specialization defect in the first Frobenius band $p\le d<2p$, with exact boundary rules and a separate plan for $p=2$.

\textbf{Route.} Use the proved high-coefficient recurrence to remove the settled sectors. Concentrate on the remaining low-coefficient binary-hyperplane tensor comparison, then track the actual elementary source through the required Witt changes. The order-five obstruction must be accounted for explicitly; a universal relative-saturation lemma is no longer a viable shortcut.

\textbf{Success.} A formula valid for all $m$ in its stated prime and degree range, with witnesses and independent torsion detectors. A further table of ranks would not meet this criterion. If a full band is too ambitious, a theorem characterizing which obstruction modules control its factors would still be a meaningful general advance.

\subsection{Priority four: structural exceptional and multiblock arithmetic}
Two worthwhile questions remain distinct. The first is a root-theoretic explanation of the complete E6/E7 relative formulas: why the observed exponents and first-gap degrees occur, and when an analogous pattern holds for other root data. The second is the ambient quotient $H/N$, especially where root and weight lattices agree. A theorem about one does not solve the other.

For classical multiblock flags, the natural target is the complete ambient B/D joint-image filtration. Start from the established primitive separate-image sum inside the joint algebra, then identify the missing ambient mixed-product lattice. The right proof should be compatible with many blocks or explain the exact obstruction to such compatibility. It should not simply re-run the single-tail Smith routine for a longer list of examples.

\subsection{Parallel review and longer-term geometry}
The completed adjoint fiber theorem should undergo focused external review before a uniform replacement proof is made the dominant project. The review has concrete tasks: audit the scheme-preserving reduction, the E8 characteristic-zero certificate implication, and the passage through central covers. Independent review means a mathematician's assessment of the argument, not merely a second run of the same software.

A future uniform proof needs an additional identity forcing the singular image to be nondominant. The radial Casimir identity, finite central-image statements, and Hessian symmetry alone have not supplied it. Positive excess in E6 and persistent nilpotents for altered characters are mandatory controls. A useful new lemma must imply a scheme-theoretic conclusion beyond those already-known constraints.

Degree questions should similarly be organized around an effective boundary theorem for a family of characters, with the correct Newton lattice and all normal multiplicities. Minimum Cayley degree, target branching, and nonproperness are separate projects. Historical claims of openness should be checked anew before becoming a proposal's central motivation.

\section{Connections, limitations, and standards of evidence}
\subsection{How the work extends the papers}
The connection is direct at the level of maps. The source is obtained using the same $\Theta_\rho^*$, the cohomological map is the same $\xi_\rho^P$, and the normalized classical generators are those of the published calculations. The extension changes the questions asked of these maps:
\begin{enumerate}
\item The explicit weight formula becomes a general arithmetic description of the full integral torus source.
\item Complex or stable injectivity becomes a finite-rank study of kernels, saturation, and every Smith multiplicity for selected classical sources.
\item Abstract divisibility becomes explicit divided classes and integral multiplication in the saturation.
\item The same integral polynomial characters enter K-theory, where support filtrations and equivariant specialization reveal additional arithmetic.
\item Before the Borel quotient, the same trace map supplies quotient coordinates, differential factors, degree corrections, and actual-fiber equations.
\end{enumerate}
The ambient Borel, BGG, Schubert, characteristic-map, and harmonic theories remain published foundations. This program does not claim to have introduced them. The candidate contributions lie in the precise representation-dependent image and fiber statements, subject to the proof and novelty qualifications already stated.

\subsection{Hiller, folding, and adjacent geometry}
Hiller's chain/intersection approach suggests ways to interpret products and incidence filtrations in Schubert calculus \cite{Hiller}. A comparison with the Smith prefixes would need an integral map and a primitivity theorem. Matching a chain count to a multiplicity polynomial would not establish the integral identification.

Diagram folding gives a different connection. For a diagram-stable representation, trace-orthogonal projection is compatible with restriction to the connected fixed subgroup. The full integral source can nevertheless enlarge after restriction. The recorded SL/Sp, E6/F4, and Spin8/G2 comparisons exhibit this defect. A folding explanation of an E7 multiplicity expression therefore requires more than noticing a smaller-type Weyl polynomial: it must identify the image lattice and account for its index.

The square-zero varieties suggest Springer-fiber and Spaltenstein geometry, but the incidence conditions must be matched exactly. This is a promising bridge between classical representation combinatorics and the denominator filtration. The most useful next connection would prove an integral restriction-kernel statement, not merely provide another interpretation of the same rational generating function.

\subsection{What remains unproved}
The principal unresolved claims are: the higher-block intermediate kernel equality; full type A specialization multiplicities beyond the stated ranges; ambient B/D joint-image factors; general ambient exceptional image formulas; a uniform root explanation of the completed adjoint case theorem; and broad exceptional nonemptiness/excess classifications. General direct-sum lattice formulas, bad-prime invariant sources, and arbitrary-representation boundary degrees remain further structural questions.

Several stronger statements have already failed and should not reappear as tacit assumptions. Images need not commute faithfully with specialization. The torus lattice alone need not determine the K-image. Universal zero transverse excess is false. Global regularity preservation is much rarer than generic fiber smoothness. A naive root-power extension of the square-generator ideal is false. Rational Hilbert-series equality does not prove integral saturation, and a reduced-support calculation does not detect embedded nilpotents.

\subsection{Verification policy}
Each claimed theorem should have a written argument, a source attribution, and a sharply delimited computational dependency. Independent checks are strongest when they change the mathematical route: BGG versus Chevalley for Schubert coefficients; symmetric-function versus native integer straightening for lattices; companion matrices versus Newton identities for character formulas. Agreement between two wrappers around the same engine is not independent verification of that engine.

The latest ordinary package records 96 B embeddings, 124 constructed D components, 28 actual D embeddings, and seven D product checks; separate native Chevalley calculations cover 26 B and 28 D components. The two-variable audit verifies monic integral reductions, including 283 S-pairs over fifteen parameter choices. These checks detect implementation and transcription errors. The all-rank claims rest on the written filtration and recurrence proofs. The statement's rendering and link checks are a separate editorial audit and do not strengthen the mathematical proof status.

\status{The broader framework and exact failed generalizations are indexed in \record{RESEARCH_LEDGER.md}{the ledger}. Folding is treated in \record{OTHER_THEORIES_FOLDING.md}{the restriction note}. The release record \record{classical-adapted-basis-release.json}{binds the latest ordinary checks} to their scripts and proof inputs.}

\appendix
\section{Result inventory and review guide}
This inventory records the principal results, rather than every historical lemma or finite calculation. ``Proved'' in the status column means a written repository proof with the stated internal checks, unless the row says published. It does not establish literature priority. The linked baseline fixes what a collaborator is being asked to review.

\small
\begin{longtable}{>{\raggedright\arraybackslash}p{0.13\textwidth}>{\raggedright\arraybackslash}p{0.38\textwidth}>{\raggedright\arraybackslash}p{0.39\textwidth}}
\toprule
Record & Result and scope & Status and remaining qualification\\\midrule\endhead
BASE-02/03 & Explicit Springer formula and polynomial-character Borel maps & Published Rogers/Kumar framework; Parts II/III supply stable and explicit families.\\
SRC-01/03 & Full torus source is $\Sym_\Z\Lambda_\rho$; Levi source is its invariant ring & Proved in the trace-nondegenerate simply connected simple scope; Levi invariants need not be generated in degree one.\\
SRC-02 & $\Lambda_\lambda=I_\lambda(Q+\Z\lambda)^*$ & Proved for every nontrivial irreducible finite simple type; direct sums require combined coefficient lattices.\\
IMG-10/11 & Torsion-index bound and exact defining projective factors & Proved full literal source statements; general bad-prime multiplicities not determined.\\
IMG-02--05 & C primitivity; B factors; D square equality and full gluing & Proved all-rank normalized single-tail results; literal scaling must be specified separately.\\
IMG-08/09 & C joint primitivity and B/D relative sum theorem & Proved for specified two-step generators; ambient B/D joint factors remain open.\\
IMG-12/20 & E6 all 63 and E7 all 127 proper relative root/weight comparisons & Complete, partly computer-assisted; not complete ambient exceptional images.\\
IMG-44/45 & Adapted B/D bases and full D multiplication & Proved constructions from previous prefix/gluing theorems; independently checked finite embeddings.\\
IMG-46/47 & Geometric denominator formula; top layer; all layers for $\ell=2$ & Proved; $J_e=K_e$ for intermediate $e$ and $\ell\ge3$ remains conjectural.\\
IMG-27/31 & Filtered finite K-image comparison; symplectic projective Smith list & Proved in the specified literal-source scope; not a lattice-only K-image classification.\\
IMG-37/39 & Type A K-specialization faithfulness and exact exponent & Proved all $m$: $\Delta_2=0$, exponent $m$ for $m\ge3$.\\
IMG-40--43 & Partition factors and corrections through $p+3$ & Proved only in the prime/degree ranges stated in Section 12; full band and higher 2-primary pieces open.\\
IMG-32/33 & Specified projective quantum and evaluation/level affine images & Proved benchmark constructions; arbitrary partial-flag lifts and affine descent remain separate.\\
GEO-05 and degree notes & E6/E7 minimal/adjoint generic degrees and further families & Exact working deductions with boundary dependencies; not minimum Cayley degrees.\\
GEO-12/13 & Symplectic sharp harmonic localization and regularity classification & Written internally audited deductions; external review recommended.\\
GEO-90--92 & All-simple proper-Levi adjoint generic regular-source smoothness & Completed by cases; E8 certificate has an explicit engine dependency; uniform explanation open.\\
GEO-94--97 & Semisimple extension; polynomial-character survival; classical blocks & Proved with exact hypotheses; arbitrary representations and exceptional excess classification open.\\
\bottomrule
\end{longtable}
\normalsize

\subsection*{Suggested reading paths}
For a collaborator primarily interested in Parts II and III, read Sections 2--9, then the B/D proof notes and the two-variable straightening proof. For a collaborator interested in specialization, read Sections 4--5 and 10--13 before the recurrence and Witt-comparison notes. For a geometric reviewer of the adjoint theorem, read Sections 14--17 and then the isolated review package, keeping the common scheme reduction distinct from the exceptional certificates.

An initial external review should focus on either the primitive-prefix/gluing chain or the actual-fiber/E8 chain. It can identify gaps, simplify arguments, or recognize prior results without endorsing the entire program.
